\section{Detailed Planning}
\label{app:planning}

This appendix expands the schedule of Section~\ref{sec:planning} into the weekly rhythm that
produced it, and records the two planning revisions made during the mission.

\subsection*{Weekly cadence}

Table~\ref{tab:cadence} lists the standing meetings that structured every week of the
mission.

\begin{table}[htbp]
\centering
\caption{Standing weekly meetings}
\label{tab:cadence}
\small
\begin{tabularx}{\linewidth}{>{\raggedright\arraybackslash}p{2.2cm} >{\raggedright\arraybackslash}p{3.4cm} >{\raggedright\arraybackslash}p{3.2cm} X}
\toprule
\textbf{Day} & \textbf{Meeting} & \textbf{With} & \textbf{Purpose} \\
\midrule
Monday & Planning, Green Supply Chain activity & \'E.\ Lardy & Set the objective for the
week \\
Tuesday & Scrum & \'E.\ Lardy & Progress, blockages \\
Wednesday & Scrum & \'E.\ Lardy & Progress, blockages \\
Wednesday & Supply chain regular point & S.\ Perthuisot & Scientific direction and
arbitration \\
Thursday & Scrum & \'E.\ Lardy & Progress, blockages \\
Friday & Demonstration & \'E.\ Lardy, S.\ Perthuisot & Show working output, not describe it
\\
\bottomrule
\end{tabularx}
\end{table}

Six standing contact points a week is a demanding cadence for a research placement. Its
effect was to bound how long any misconception could survive: the Friday demonstration in
particular required something to be shown rather than reported, which repeatedly exposed
gaps between what was believed to work and what did.

\subsection*{Planning revisions}

\paragraph{June: validation over extension.}
The original plan allocated the second half of the mission to widening the model's scope. It
was reallocated to validating what already existed. The argument accepted at the regular
point was that extending an unvalidated model compounds risk rather than reducing it. This
revision is why the mission produced a measured result rather than a larger unmeasured
system.

\paragraph{July: substitution of the campaign population.}
The campaign was designed for eight human consultants across nine working days. When that
could not be secured within the mission window, the protocol was executed with isolated
agents instead. The record format was deliberately identical, so the substitution changed the
population without changing the instrument. Every behavioural conclusion was weakened
accordingly in the reporting rather than being stated at the strength the original design
would have supported.

\subsection*{Comparison of planned against actual}

Table~\ref{tab:planning_ecart} compares the time allowed for each phase with the time it
actually took.

\begin{table}[htbp]
\centering
\caption{Planned against actual, by phase}
\label{tab:planning_ecart}
\small
\begin{tabularx}{\linewidth}{>{\raggedright\arraybackslash}p{3.6cm} >{\centering\arraybackslash}p{2.2cm} >{\centering\arraybackslash}p{2.2cm} X}
\toprule
\textbf{Phase} & \textbf{Planned} & \textbf{Actual} & \textbf{Comment} \\
\midrule
Framing & 3 weeks & 2.5 weeks & On plan \\
First urgency model and pipeline & 6 weeks & 5.5 weeks & Stopped early once judged
inadequate \\
Capacity from imagery & 4 weeks & 8 weeks & Ran in parallel; the original objective proved
unreachable and the problem had to be reframed \\
SupplyScore build & 4 weeks & 2 weeks & Interfaces already established by the abandoned
pipeline \\
Experimental design & 2 weeks & 3 weeks & Freezing the protocol properly took longer than
allowed \\
First campaign and diagnosis & 3 weeks & 3 weeks & On plan \\
Independent validation & \emph{unplanned} & 3.5 weeks & Added after the diagnosis:
repair, second
chain, baseline comparison \\
Calibration campaigns & 4 weeks & \emph{planned} & 25 August to 19 September \\
Vision model & 3 weeks & \emph{planned} & 7 to 26 September, overlapping the above \\
Scientific report and defence & 3 weeks & \emph{planned} & 28 September to 16 October \\
\bottomrule
\end{tabularx}
\end{table}

The overall schedule held. Two overruns had to be absorbed. The imagery workstream took twice
its allowance because the original objective turned out to be physically unreachable and the
problem had to be reframed before anything could be built; the time came out of the
SupplyScore build, which the abandoned pipeline had already de-risked. The validation phase
was unplanned entirely and consumed three and a half weeks, released by the framing and build
phases finishing early. The one systematic underestimate was the effort required to freeze an
experimental protocol to a standard where the hypotheses could not later be adjusted, which
took half again the time allowed and was worth every day.

\subsection*{What the remaining two months carry}

The mission runs to 16 October and this report is submitted on 24 August. The remaining time
carries no new construction, and it is loaded to roughly the same weekly effort as the phases
already closed.

\begin{enumerate}
    \item \emph{Calibration campaigns, 25 August to 19 September.} New serious-game chains,
    built the way the airframer chain was, run to widen the evidence base. This is the single
    action that moves every sample-bound figure in Section~\ref{sec:conclusion_limites}, and
    it is also what allows the forecast layer to be recalibrated on more than one scenario.
    \item \emph{Vision model, 7 to 26 September, overlapping the above.} The facade
    segmentation model of Section~\ref{sec:realisation_volume} carries a global bias of
    $-35\%$ on the derived dock count. Annotated volume, not architecture, drove every gain
    measured so far, so the work is data collection and re-training rather than redesign.
    \item \emph{Scientific report and defence, 28 September to 16 October.} The
    \emph{Bilan Scientifique et Technique} that capitalises the whole internship for the
    programme, and the defence.
\end{enumerate}

Any rate quoted in this report may therefore move. The constructions and the diagnoses will
not.
